\section{Positive observation models and history reduction}
\label{sec:v10-positive-models}

We give explicit model consequences in which the marked law,
continuity and contraction hypotheses can be checked from primitive
data. These results concern the models specified here. They do not
identify a finite-state model with an unconstructed microscopic
billiard or kinetic model.

Let $Z_t\in[n]$, $a\in A$ with $A$ finite, and suppose
\begin{equation}\label{eq:v10-positive-data}
 \epsilon\le P_a(z,z')\le M,\qquad
 Y_{t+1}=m_{a,Z_{t+1}}+\sigma_{a,Z_{t+1}}\xi_{t+1},
 \qquad \xi_t\sim N(0,1),
\end{equation}
where $0<\epsilon<M$, $|m_{a,z}|\le m_*$ and
$0<\sigma_-\le\sigma_{a,z}\le\sigma_+$. All primitive coefficients
may depend continuously on a compact unknown parameter, with the
same bounds. The initial hidden probabilities and any finite marking
kernel also depend continuously on this parameter. The observation is of the \emph{new} state. Its joint
transition/report kernel is
\[
 K_a(z',dy\mid z)=P_a(z,z')g_{a,z'}(y)\,dy.
\]
Finite terminal or path marks are permitted.

\begin{theorem}[A continuous-report consumer]\label{thm:v10-hmm}
For every fixed horizon, the instruments
\eqref{eq:v10-positive-data} satisfy the compact dominated hypotheses
of Theorem~\ref{thm:v10-compact} for bounded losses continuous in the
primitive parameter. Their common-encoder minimax problems have
attaining private designs and attaining shared mixtures.

Partition $[-R,R]$ into intervals of length at most $h$, where $R>m_*$,
and use two additional tail cells. Reconstruction is uniform on
interior intervals and uses any fixed probability density on each
tail. For every memory width and every admissible controller, the
resulting marked finite-report approximation has risk error at most
$T e(h,R)$, where
\begin{equation}\label{eq:v10-bin-error}
 e(h,R)=\frac{h}{2\sigma_-}\sqrt{\frac2\pi}
       +2\exp\left(-\frac{(R-m_*)^2}{2\sigma_+^2}\right).
\end{equation}
The corresponding excess-risk error is at most $2T e(h,R)$ when
baselines are recomputed.
\end{theorem}
\begin{proof}
Use a positive Gaussian probability as the reference measure on each
report space and counting probability on finite marks. With an
open-loop action word fixed, the joint marked density is a finite
sum over hidden paths of products of transition probabilities,
Gaussian densities and terminal marking probabilities. Gaussian
location-scale densities depend continuously in $L^1$ on their
parameters when the variance is bounded away from zero. This follows,
for example, by pointwise convergence and equality of their integrals,
or directly by truncating uniformly Gaussian tails and using uniform
continuity on a compact interval. Products of probability densities
are continuous in $L^1$ by telescoping, and so are finite mixtures.
Compactness of the primitive parameter set gives compactness of the
path-density family in $L^1$. Multiplication by bounded continuously
parameterized losses preserves the required effect continuity.
Theorem~\ref{thm:v10-compact} applies.

If $I$ is an interval of length at most $h$ and $\bar g_I$ is its
average density, the fundamental theorem of calculus gives
\[
 \int_I|g-\bar g_I|\le h\int_I|g'|.
\]
Indeed average $|g(x)-g(y)|$ over $y\in I$ and bound the difference
by the integral of $|g'|$ between $x,y$. For a Gaussian density,
$\int|g'|=\sqrt{2/\pi}/\sigma$. On the two tails, replacing a
conditional density changes total variation by at most their total
probability, bounded by the exponential term in
\eqref{eq:v10-bin-error}. Thus each conditional observation channel
has total variation error at most $e(h,R)$ uniformly in hidden state,
control and parameter. Couple the hidden transition identically and
maximally couple the original and reconstructed observations. As
long as observations and random tapes have matched, both controllers
make the same decisions. The probability of a first mismatch during
$T$ steps is at most $Te(h,R)$. On its complement the entire hidden
path, all reports and the actual mark agree. Bounded losses change
by at most this probability. The report conversions are transient
operations, so the memory width is unchanged. Comparing the
unrestricted baselines in the same way finishes the proof.
\end{proof}

\subsection{A proved forgetting estimate}
For positive probability rows $u,v$, put
\[
 d_H(u,v)=\max_i\log(u_i/v_i)-\min_i\log(u_i/v_i).
\]
Multiplying both rows by the same positive diagonal matrix and then
normalizing leaves $d_H$ unchanged. We use the elementary positive
matrix contraction in its explicit form.

\begin{lemma}[Positive-matrix contraction]\label{lem:v10-hilbert}
If the entries of a matrix $P$ belong to $[\epsilon,M]$, then
\[
 d_H(uP,vP)\le\rho d_H(u,v),\qquad
 \rho=\frac{M-\epsilon}{M+\epsilon}<1.
\]
For arbitrary nonnegative probability rows, the first images satisfy
$d_H(uP,vP)\le2\log(M/\epsilon)$.
\end{lemma}
\begin{proof}
Set $x_i=\log(u_i/v_i)$ and interpolate $v_i\exp(sx_i)$,
$0\le s\le1$. The derivative of the logarithm of output coordinate
$j$ is the expectation of $x_i$ under weights proportional to
$v_i\exp(sx_i)P_{ij}$. For two output coordinates $j,k$, the likelihood
ratio of their normalized weight distributions has maximum divided
by minimum at most $(M/\epsilon)^2$. If two probability distributions
have likelihood ratios in an interval $[a,b]$ with $b/a\le C$, their
total variation is at most $(\sqrt C-1)/(\sqrt C+1)$.
To verify this bound, put all ratios below one at the lower endpoint
and all ratios above one at the upper endpoint; the mass constraint
makes the resulting variation at most
$(1-a)(b-1)/(b-a)$. Maximizing over $b/a\le C$, $a\le1\le b$,
gives $a=C^{-1/2}$, $b=C^{1/2}$ and the displayed value.
Differences of the two derivative expectations are therefore at
most $\rho\osc(x)$. Integrate in $s$ and maximize over $j,k$.
At $s=0$ all output log ratios vanish, proving the first claim.
For the second, every coordinate of either probability row times
$P$ lies in $[\epsilon,M]$; its coordinatewise ratio lies in
$[\epsilon/M,M/\epsilon]$.
\end{proof}

\begin{theorem}[Uniform finite-history recovery]\label{thm:v10-forgetting}
Fix the primitive parameter. Two filters for
\eqref{eq:v10-positive-data}, started with arbitrary initial
probabilities and fed the same actions and observations for $L\ge1$
steps, satisfy
\[
 \TV(\pi_L,\widetilde\pi_L)
 \le \frac12\log(M/\epsilon)\rho^{L-1}.
\]
In particular a last-$L$-history filter, initialized with a fixed
probability at its left endpoint, approximates the full filter with
this error uniformly in the earlier history. With $N$ report cells,
a literal suffix implementation uses at most $(|A|N)^L$ labels,
besides a finite beginning-of-stream convention. For any $[0,1]$-valued terminal loss on the current hidden state,
an optimal suffix-based action has risk at
most twice this filtering error above the full-filter optimum.
\end{theorem}
\begin{proof}
Prediction multiplies a posterior row by $P_a$. Conditioning on the
current report multiplies it by the positive diagonal matrix with
entries $g_{a,z'}(y)$ and normalizes. The latter operation preserves
$d_H$, and Lemma~\ref{lem:v10-hilbert} contracts the former. The first
prediction gives diameter at most $2\log(M/\epsilon)$ and the next
$L-1$ predictions give the factor $\rho^{L-1}$. For probability rows,
$\TV(u,v)\le\tanh(d_H(u,v)/4)\le d_H(u,v)/4$; the first inequality
follows from the same two-endpoint likelihood-ratio bound in the
lemma. Restart the truncated filter with the fixed prior at time
$t-L$ and apply this result to obtain the suffix estimate.
Its raw suffix is updated by shifting and appending the current
control/report cell. Finally every $[0,1]$ loss expectation changes
by at most the total variation of the two posteriors. Compare the
optimal action under each posterior through the other posterior,
which costs at most twice that distance.
\end{proof}

This contraction is not asserted for arbitrary past-path target marks:
augmenting the hidden state to retain such a mark can destroy strict
positivity. The general marked duality retains these marks without
using the forgetting assertion.

The parameter must be fixed in this forgetting assertion. A static
unknown phase need not be forgotten, and applying the assertion to
its posterior would be false. The compact robust theorem applies to
the unknown-parameter experiment without that extra assertion. Also,
summable variations alone have not been used to infer exponential
decay: here it follows from the explicit positive-matrix contraction.
The distinction is important for history-kernel applications.

\section{Nonlinear finite-state semigroups}
\label{sec:v10-nisio}

Let $Q^a$ be conservative generators on $[n]$, $a\in A$ finite, and
let $v_i^a$ be bounded potentials. Define on $\R^n$
\begin{equation}\label{eq:v10-nisio}
 H_i(u)=\max_{a\in A}
 \left\{v_i^a+\sum_{j\ne i}q_{ij}^a
                 \bigl(e^{u_j-u_i}-1\bigr)\right\}.
\end{equation}
This is a nonlinear Hamiltonian. The following proof does not apply
a linear resolvent identity to it.

\begin{theorem}[Global semigroup and nonlinear resolvent]
\label{thm:v10-nisio}
The equation $\dot u=H(u)$ has a global solution for every initial
condition and defines a semigroup $S_t$ that is order preserving,
additively homogeneous and nonexpansive in the sup norm. If
$V=\max_{i,a}|v_i^a|$, then
$\norm{S_tu}_\infty\le\norm u_\infty+tV$.
For every $\lambda>0$ the equation
\[
 u-\lambda H(u)=f
\]
has a unique solution $J_\lambda f$. The resolvent is order preserving,
additively homogeneous and nonexpansive. In particular $H$ is
maximal dissipative with this sign convention.
\end{theorem}
\begin{proof}
The map $H$ is locally Lipschitz as the maximum of finitely many
smooth functions. If $u_i-v_i=\max_j(u_j-v_j)$, then
$u_j-u_i\le v_j-v_i$ for every $j$, and hence $H_i(u)\le H_i(v)$.
The upper right derivative of the maximum coordinate of the
difference of two solutions is therefore nonpositive. Applying the
same argument to the reverse difference proves sup-norm contraction
and comparison. Since $H(u+c\1)=H(u)$, uniqueness gives additive
homogeneity. Comparing with the time-dependent constant barriers
$(\max_i u_i+tV)\1$ and $(\min_i u_i-tV)\1$ gives the norm bound.
It prevents finite-time escape, so the locally defined solution is
global; uniqueness gives the semigroup law.

For the resolvent, consider the damped equation
$\dot u=f-u+\lambda H(u)$. The interval of constant barriers
\[
 (\min_i f_i-\lambda V)\1\le u\le
 (\max_i f_i+\lambda V)\1
\]
is invariant. At a maximal difference between two solutions the
upper derivative is at most the negative of that difference.
Thus its flow contracts distances by $e^{-t}$. At time one this is
a strict contraction of a complete invariant rectangle into itself,
so it has a unique fixed point. Commutation with the flow shows
that the fixed point is stationary: every later image is another
fixed point of the time-one map. It therefore solves the resolvent
equation. The same maximal-coordinate comparison applied to two
resolvent equations gives
$\max_i(u_i-\widetilde u_i)\le\max_i(f_i-\widetilde f_i)$,
which proves uniqueness, comparison and nonexpansiveness.
Translation by constants proves homogeneity.
\end{proof}

\begin{theorem}[Marked mesh approximation and stability]
\label{thm:v10-nisio-mesh}
Define the log-moment Bellman step
\[
 (F_\Delta u)_i
 =\max_a\log\left( e^{\Delta(Q^a+\diag v^a)}e^u\right)_i.
\]
For fixed $t$ and bounded initial $u$, with $\Delta=t/m$,
$F_\Delta^m u\to S_tu$ with error $O(\Delta)$, uniformly on bounded
sets. Each mesh problem, after acquisition of the initial state, has
finite report and control spaces and a bounded actual path-reward mark.
Its bounded moment effects therefore have the common-encoder
risk-body duality of Corollary~\ref{cor:v10-effects}, for compact
continuously parameterized primitive data. The seed average is taken on these
moments before taking a logarithm.

If $\widetilde Q^a,\widetilde v^a$ are another such family and
\[
 \delta_Q=\max_{i,a}\sum_{j\ne i}|q_{ij}^a-\widetilde q_{ij}^a|,
 \qquad \delta_v=\max_{i,a}|v_i^a-\widetilde v_i^a|,
\]
then on a common oscillation bound $B$ for the trajectories,
\[
 \norm{S_tu-\widetilde S_tu}_\infty
 \le t\{\delta_v+(e^B+1)\delta_Q\}.
\]
For the resolvents the same bound holds with $t$ replaced by
$\lambda$ and $B$ bounding their solutions.
\end{theorem}
\begin{proof}
The matrix exponential is a positive Feynman--Kac operator. On a
bounded set of $u$, its Taylor expansion and the logarithm give
$F_\Delta u=u+\Delta H(u)+O(\Delta^2)$, uniformly for the finite
family of generators. Positivity makes $F_\Delta$ order preserving,
and $F_\Delta(u+c\1)=F_\Delta u+c\1$. These two properties imply
sup-norm nonexpansiveness. Its iterates and the exact flow remain
in the same bounded region up to time $t$, by the constant barriers.
The local flow expansion has the same first-order term and a uniform
$O(\Delta^2)$ remainder. Telescope the errors over $m$ steps using
nonexpansiveness; the bound is $mO(\Delta^2)=O(t\Delta)$.

For a fixed control during one interval, the positive kernel
$e^{\Delta(Q^a+\diag v^a)}$ is the expectation of the exponential
integrated reward with the terminal state indicator. Keeping the
actual bounded reward integral as a mark, rather than interpreting
this nonstochastic matrix as a transition probability, gives its
proper probabilistic representation. Concatenation gives the mesh
instrument. Finite sampled states and controls have a finite report
representation; continuous reward marks can be integrated into its
bounded effect functions. For this reward-maximizing problem, use a decreasing affine
normalization of the bounded moment as a loss in $[0,1]$ when applying
the risk theorem. This reverses the optimization in the required way. For finite state/action
paths this integration needs no assumption that reward marks have a
common density: Corollary~\ref{cor:v10-effects} applies directly.
For finite parameter and task sets, one may also augment the finite
hidden state by the sampled state/control path and replace the moment
loss by its conditional expectation given that path. This gives a
finite instrument with exactly the same moment risks, to which
Theorem~\ref{thm:v9-minimax} applies. It is an equivalence for these
moment effects, not a finite-state realization of every continuous
reward mark. Continuity of the effects on a compact primitive family
follows from continuity of the finite matrix exponentials and products.

On $\osc(u)\le B$, direct subtraction of \eqref{eq:v10-nisio} gives
$\norm{H(u)-\widetilde H(u)}_\infty
\le\delta_v+(e^B+1)\delta_Q$. The maximal-coordinate comparison used
above now bounds the upper derivative of the solution difference by
this constant. Integration proves the flow estimate. Subtracting
the two resolvent equations at a maximizing coordinate proves its
resolvent counterpart. The logarithm of an average exponential
moment is not an average of logarithms; none of these steps changes
that ordering.
\end{proof}

The semigroup construction is a finite-state instance of the
logarithmic control-semigroup method of Nisio~\cite{Nisio1976} and risk-sensitive
control. Its role here is an explicit nonlinear consumer with a
proved range condition and a marked finite-mesh interface, not a
claim that finite-dimensional comparison settles the domain and
coercivity questions of an infinite-dimensional kinetic generator.

\section{Hidden phases and risk-sensitive continuation}
\label{sec:v10-phases}

Let $C$ be a finite latent phase chosen once with positive initial
weights, and let $Z_t$ follow a controlled hidden model conditional
on $C$. Its joint transition/report densities are
$K_{c,a}(z',y\mid z)$ with respect to a probability measure $\mu(dy)$.
Let $r_{c,a}(z,z')$ and the terminal reward $h_c(z)$ be bounded, and
fix $\gamma>0$. One controller must choose each action without
observing $C$. The objective in this section is to minimize the
positive exponential moment. A logarithm may be applied afterwards.

The exponential information-state mechanism is classical; see
B\"auerle and Rieder~\cite{BauerleRieder2017}. We use the following
explicit unnormalized version so that no phasewise policy is hidden
in the notation. Define a nonnegative row $\eta$ on $(c,z)$ and the tilted update
\begin{equation}\label{eq:v10-tilt}
 (\eta K_a^\gamma(y))(c,z')
 =\sum_z\eta(c,z)K_{c,a}(z',y\mid z)
                         e^{\gamma r_{c,a}(z,z')}.
\end{equation}
Unlike an ordinary phase posterior, $\eta$ includes the accumulated
reward weight. Its total mass must not be discarded without a
corresponding multiplicative value factor.

\begin{theorem}[One-policy phase recursion]\label{thm:v10-phase}
For full-history controls, the minimum remaining exponential moment
has the recursion
\[
 V_T(\eta)=\sum_{c,z}\eta(c,z)e^{\gamma h_c(z)},\qquad
 V_t(\eta)=\min_{a\in A}\int
          V_{t+1}(\eta K_a^\gamma(y))\mu(dy).
\]
The value is positively homogeneous and Borel, and an optimal
measurable selector exists. The minimum is outside the phase sum:
replacing it by separate phasewise minima is in general incorrect.
For finite-register common encoders, the exact moment risk body is
instead obtained from the coupled design of
Theorem~\ref{thm:v10-compact} (or its finite specialization), not by
providing $\eta$ as free online memory. The model satisfies that
theorem whenever its marked path densities and bounded effects
satisfy the stated domination and continuity hypotheses; in
particular finite phases with positive Gaussian observations and
bounded finite-path rewards do so for fixed primitives, or for compact
continuous primitive families as in Theorem~\ref{thm:v10-hmm}, using
Corollary~\ref{cor:v10-effects} when only effect domination is needed.
\end{theorem}
\begin{proof}
For a fixed report/action history, $\eta_t(c,z)$ is the joint density
of phase, current state and that report history, multiplied by the
exponential accumulated reward and summed over past hidden states.
This interpretation is true initially and propagates exactly by
\eqref{eq:v10-tilt}. At the terminal time summing the remaining
exponential terminal factor gives $V_T$. At time $t$ the action is
one measurable choice from the current history; it cannot depend on
an unobserved phase. For a fixed choice $a$, integrate the remaining
value of the updated density against the reference observation law.
Backward induction gives the recursion. The minimum of finitely
many Borel integrals is Borel and the least-index minimizer is a
measurable selector. Induction also gives positive homogeneity.
This is the unnormalized information-state form of risk-sensitive
dynamic programming; it accounts for the probability and reward
weights together.

For a finite-register design, fix its action and register paths and
expand the true marked path law. Its exponential reward is a bounded
positive loss after affine rescaling. The same encoder kernels appear
in every task and phase row, so the compact risk theorem applies
without giving the controller the unbounded-dimensional history or
an uncharged exact real vector. The Gaussian model's domination and
compactness follow from the finite mixture proof in
Theorem~\ref{thm:v10-hmm}. This proves the last assertion.
\end{proof}

\begin{example}[Strict failure of phasewise optimization]
\label{ex:v10-phase}
There is no observation. The equally likely phase is $C\in\{0,1\}$,
the action is $a\in\{0,1\}$, and the cost is $\1_{\{a\ne C\}}$.
The best actual exponential moment is $(1+e^\gamma)/2$, for any
mixture of the two actions. Its logarithmic value is
$\gamma^{-1}\log((1+e^\gamma)/2)$. Choosing $a=C$ separately in each
hidden phase would give moment one, but is not an admissible policy.
The tilted recursion and the common-encoder risk body both retain
the actual value.
\end{example}
